% --- Archon physics formalization source begin ---
% archon:physics
% archon:covers PhyXMiniProblems/problem_phyx_mini_0592.lean
% archon:source-report reports/phyx_mini/problem_phyx_mini_0592.source.json

\chapter{Physics problem phyx\_mini\_0592}
\label{ch:PhyXMiniProblems_problem_phyx_mini_0592}

\paragraph{Problem source.}
\#\# Physical scenario

Figure gives their spatial separation \$\textbackslash{}Delta x\$ according to the \$S\$ observer as a function of \$\textbackslash{}Delta t\$ for a range of \$\textbackslash{}Delta t\$ values. The vertical axis scale is set by \$\textbackslash{}Delta x\_a = 10.0 \textbackslash{}, \textbackslash{}text\{m\}\$.

\#\# Question

What is \$\textbackslash{}Delta x'\$?

\#\# Answer choices

- A: A: 2.0m

- B: B: 1.3m

- C: C: 0.79m

- D: D: 3.0m

\#\# Figure caption (auxiliary; use the image as primary evidence)

The image is a graph with a labeled coordinate system. The x-axis is labeled as "Δt' (ns)," representing time in nanoseconds, while the y-axis is labeled "Δx (m)," representing distance in meters. A straight line is plotted diagonally from the origin upwards to the right, indicating a positive correlation between Δt' and Δx. Additionally, there is a vertical label to the left, "Δxₐ," positioned above the graph. The line is dark and bold against a grid background with increments at 0, 4, and 8 on the x-axis.

\#\# Dataset metadata

Relativity; Spatial Relation Reasoning, Physical Model Grounding Reasoning

\paragraph{Recorded answer/context.}
C: C: 0.79m

\paragraph{Figure/image path.}
/root/proposal\_for\_physic/science-mango/phyx\_mini\_run/phyx\_data/test\_image/592.png

\paragraph{Formalization target.}
create a compiling Lean file with sorry bodies at `PhyXMiniProblems/problem\_phyx\_mini\_0592.lean`. The Lean declarations must preserve the physical quantities, dimensions or dimensional roles, figure labels, governing-law hypotheses, and final relation expressed by this problem.
Use Mathlib/Physlib names found through LeanExplore where available. If a physics API is missing, introduce faithful local abstractions rather than scalar placeholder aliases.

\begin{theorem}[Physics formalization target]
\label{thm:physics:phyx_mini_0592:target}
\lean{PhyXMiniProblems.ProblemPhyXMini0592.primedSeparation_matches_choice_C}
\uses{def:physics:phyx-mini-0592:phyxminiproblems-problemphyxmini0592-separationinmeters, def:physics:phyx-mini-0592:phyxminiproblems-problemphyxmini0592-vacuumlightspeedinmeterspernanosecond, def:physics:phyx-mini-0592:phyxminiproblems-problemphyxmini0592-relativisticseparationsetup, def:physics:phyx-mini-0592:phyxminiproblems-problemphyxmini0592-matchesrelativisticseparationscenario, def:physics:phyx-mini-0592:phyxminiproblems-problemphyxmini0592-matchessuppliedseparationtimegraph, def:physics:phyx-mini-0592:phyxminiproblems-problemphyxmini0592-hasphysicalrelativisticparameters, def:physics:phyx-mini-0592:phyxminiproblems-problemphyxmini0592-satisfiesinverselorentzspatialtransformation, def:physics:phyx-mini-0592:phyxminiproblems-problemphyxmini0592-matchesdisplayedprecision, def:physics:phyx-mini-0592:phyxminiproblems-problemphyxmini0592-isnearestanswerchoice}
The assigned autoformalize agent should translate this physics problem into Lean declarations in the covered file, with theorem and lemma proof bodies written as `by sorry`.
\end{theorem}
\begin{proof}
This is an autoformalization task, not a proof task. Produce statements that can later be proved by the physics prover without weakening the physical model.
\end{proof}

% --- Archon declaration topology begin ---
\section{Declaration topology}
\begin{definition}[Spatial Separation]
  \label{def:physics:phyx-mini-0592:phyxminiproblems-problemphyxmini0592-spatialseparation}
  \lean{PhyXMiniProblems.ProblemPhyXMini0592.SpatialSeparation}
  A signed, unit-independent physical spatial separation
\end{definition}
\begin{proof}
  This is the declared model definition of Spatial Separation.
\end{proof}
\begin{definition}[Time Interval]
  \label{def:physics:phyx-mini-0592:phyxminiproblems-problemphyxmini0592-timeinterval}
  \lean{PhyXMiniProblems.ProblemPhyXMini0592.TimeInterval}
  A signed, unit-independent physical time interval
\end{definition}
\begin{proof}
  This is the declared model definition of Time Interval.
\end{proof}
\begin{definition}[Speed Magnitude]
  \label{def:physics:phyx-mini-0592:phyxminiproblems-problemphyxmini0592-speedmagnitude}
  \lean{PhyXMiniProblems.ProblemPhyXMini0592.SpeedMagnitude}
  A nonnegative, unit-independent relative-speed magnitude
\end{definition}
\begin{proof}
  This is the declared model definition of Speed Magnitude.
\end{proof}
\begin{definition}[meter Nanosecond Units]
  \label{def:physics:phyx-mini-0592:phyxminiproblems-problemphyxmini0592-meternanosecondunits}
  \lean{PhyXMiniProblems.ProblemPhyXMini0592.meterNanosecondUnits}
  Coherent units with metres for length and nanoseconds for time
\end{definition}
\begin{proof}
  This is the declared model definition of meter Nanosecond Units.
\end{proof}
\begin{definition}[separation In Meters]
  \label{def:physics:phyx-mini-0592:phyxminiproblems-problemphyxmini0592-separationinmeters}
  \lean{PhyXMiniProblems.ProblemPhyXMini0592.separationInMeters}
  \uses{def:physics:phyx-mini-0592:phyxminiproblems-problemphyxmini0592-spatialseparation}
  Read a signed spatial separation in metres
\end{definition}
\begin{proof}
  This is the declared model definition of separation In Meters.
\end{proof}
\begin{definition}[time Interval In Nanoseconds]
  \label{def:physics:phyx-mini-0592:phyxminiproblems-problemphyxmini0592-timeintervalinnanoseconds}
  \lean{PhyXMiniProblems.ProblemPhyXMini0592.timeIntervalInNanoseconds}
  \uses{def:physics:phyx-mini-0592:phyxminiproblems-problemphyxmini0592-timeinterval, def:physics:phyx-mini-0592:phyxminiproblems-problemphyxmini0592-meternanosecondunits}
  Read a signed time interval in nanoseconds
\end{definition}
\begin{proof}
  This is the declared model definition of time Interval In Nanoseconds.
\end{proof}
\begin{definition}[speed In Meters Per Nanosecond]
  \label{def:physics:phyx-mini-0592:phyxminiproblems-problemphyxmini0592-speedinmeterspernanosecond}
  \lean{PhyXMiniProblems.ProblemPhyXMini0592.speedInMetersPerNanosecond}
  \uses{def:physics:phyx-mini-0592:phyxminiproblems-problemphyxmini0592-speedmagnitude, def:physics:phyx-mini-0592:phyxminiproblems-problemphyxmini0592-meternanosecondunits}
  Read a speed magnitude in metres per nanosecond
\end{definition}
\begin{proof}
  This is the declared model definition of speed In Meters Per Nanosecond.
\end{proof}
\begin{definition}[vacuum Light Speed In Meters Per Nanosecond]
  \label{def:physics:phyx-mini-0592:phyxminiproblems-problemphyxmini0592-vacuumlightspeedinmeterspernanosecond}
  \lean{PhyXMiniProblems.ProblemPhyXMini0592.vacuumLightSpeedInMetersPerNanosecond}
  \uses{def:physics:phyx-mini-0592:phyxminiproblems-problemphyxmini0592-meternanosecondunits}
  Physlib's exact vacuum light speed, read in metres per nanosecond
\end{definition}
\begin{proof}
  This is the declared model definition of vacuum Light Speed In Meters Per Nanosecond.
\end{proof}
\begin{definition}[time Interval Of Nanoseconds]
  \label{def:physics:phyx-mini-0592:phyxminiproblems-problemphyxmini0592-timeintervalofnanoseconds}
  \lean{PhyXMiniProblems.ProblemPhyXMini0592.timeIntervalOfNanoseconds}
  \uses{def:physics:phyx-mini-0592:phyxminiproblems-problemphyxmini0592-timeinterval, def:physics:phyx-mini-0592:phyxminiproblems-problemphyxmini0592-meternanosecondunits}
  The dimensionful time interval having the specified nanosecond readout
\end{definition}
\begin{proof}
  This is the declared model definition of time Interval Of Nanoseconds.
\end{proof}
\begin{definition}[Inertial Frame Label]
  \label{def:physics:phyx-mini-0592:phyxminiproblems-problemphyxmini0592-inertialframelabel}
  \lean{PhyXMiniProblems.ProblemPhyXMini0592.InertialFrameLabel}
  The unprimed observer frame and the relatively moving primed frame
\end{definition}
\begin{proof}
  This is the declared model definition of Inertial Frame Label.
\end{proof}
\begin{definition}[Axial Direction]
  \label{def:physics:phyx-mini-0592:phyxminiproblems-problemphyxmini0592-axialdirection}
  \lean{PhyXMiniProblems.ProblemPhyXMini0592.AxialDirection}
  Direction along the common one-dimensional spatial axis
\end{definition}
\begin{proof}
  This is the declared model definition of Axial Direction.
\end{proof}
\begin{definition}[Graph Axis]
  \label{def:physics:phyx-mini-0592:phyxminiproblems-problemphyxmini0592-graphaxis}
  \lean{PhyXMiniProblems.ProblemPhyXMini0592.GraphAxis}
  The two coordinate axes in the supplied graph
\end{definition}
\begin{proof}
  This is the declared model definition of Graph Axis.
\end{proof}
\begin{definition}[Graph Axis Label]
  \label{def:physics:phyx-mini-0592:phyxminiproblems-problemphyxmini0592-graphaxislabel}
  \lean{PhyXMiniProblems.ProblemPhyXMini0592.GraphAxisLabel}
  Literal physical labels attached to the two graph axes
\end{definition}
\begin{proof}
  This is the declared model definition of Graph Axis Label.
\end{proof}
\begin{definition}[Vertical Scale Label]
  \label{def:physics:phyx-mini-0592:phyxminiproblems-problemphyxmini0592-verticalscalelabel}
  \lean{PhyXMiniProblems.ProblemPhyXMini0592.VerticalScaleLabel}
  The separate scale marker printed at the top of the vertical axis
\end{definition}
\begin{proof}
  This is the declared model definition of Vertical Scale Label.
\end{proof}
\begin{definition}[Plotted Curve Shape]
  \label{def:physics:phyx-mini-0592:phyxminiproblems-problemphyxmini0592-plottedcurveshape}
  \lean{PhyXMiniProblems.ProblemPhyXMini0592.PlottedCurveShape}
  The qualitative shape of the heavy plotted curve
\end{definition}
\begin{proof}
  This is the declared model definition of Plotted Curve Shape.
\end{proof}
\begin{definition}[Separation Time Graph Figure]
  \label{def:physics:phyx-mini-0592:phyxminiproblems-problemphyxmini0592-separationtimegraphfigure}
  \lean{PhyXMiniProblems.ProblemPhyXMini0592.SeparationTimeGraphFigure}
  \uses{def:physics:phyx-mini-0592:phyxminiproblems-problemphyxmini0592-spatialseparation, def:physics:phyx-mini-0592:phyxminiproblems-problemphyxmini0592-graphaxis, def:physics:phyx-mini-0592:phyxminiproblems-problemphyxmini0592-graphaxislabel, def:physics:phyx-mini-0592:phyxminiproblems-problemphyxmini0592-verticalscalelabel, def:physics:phyx-mini-0592:phyxminiproblems-problemphyxmini0592-plottedcurveshape}
  Semantic transcription of the graph. The scalar fields are named coordinate readouts from the bitmap; the vertical scale itself remains a dimensionful length
\end{definition}
\begin{proof}
  This is the declared model definition of Separation Time Graph Figure.
\end{proof}
\begin{definition}[displayed Line Slope Meters Per Nanosecond]
  \label{def:physics:phyx-mini-0592:phyxminiproblems-problemphyxmini0592-displayedlineslopemeterspernanosecond}
  \lean{PhyXMiniProblems.ProblemPhyXMini0592.displayedLineSlopeMetersPerNanosecond}
  \uses{def:physics:phyx-mini-0592:phyxminiproblems-problemphyxmini0592-separationtimegraphfigure}
  Slope of the displayed straight line, in metres per nanosecond
\end{definition}
\begin{proof}
  This is the declared model definition of displayed Line Slope Meters Per Nanosecond.
\end{proof}
\begin{definition}[Relativistic Separation Setup]
  \label{def:physics:phyx-mini-0592:phyxminiproblems-problemphyxmini0592-relativisticseparationsetup}
  \lean{PhyXMiniProblems.ProblemPhyXMini0592.RelativisticSeparationSetup}
  \uses{def:physics:phyx-mini-0592:phyxminiproblems-problemphyxmini0592-spatialseparation, def:physics:phyx-mini-0592:phyxminiproblems-problemphyxmini0592-timeinterval, def:physics:phyx-mini-0592:phyxminiproblems-problemphyxmini0592-speedmagnitude, def:physics:phyx-mini-0592:phyxminiproblems-problemphyxmini0592-inertialframelabel, def:physics:phyx-mini-0592:phyxminiproblems-problemphyxmini0592-axialdirection, def:physics:phyx-mini-0592:phyxminiproblems-problemphyxmini0592-separationtimegraphfigure}
  Independent physical quantities in the relativity experiment. deltaXPrime is the unknown primed-frame spatial separation. deltaXAccordingToS is the family of unprimed-frame separations whose metre readouts are plotted as the primed time interval varies
\end{definition}
\begin{proof}
  This is the declared model definition of Relativistic Separation Setup.
\end{proof}
\begin{definition}[speed Fraction Of Light]
  \label{def:physics:phyx-mini-0592:phyxminiproblems-problemphyxmini0592-speedfractionoflight}
  \lean{PhyXMiniProblems.ProblemPhyXMini0592.speedFractionOfLight}
  \uses{def:physics:phyx-mini-0592:phyxminiproblems-problemphyxmini0592-speedinmeterspernanosecond, def:physics:phyx-mini-0592:phyxminiproblems-problemphyxmini0592-vacuumlightspeedinmeterspernanosecond, def:physics:phyx-mini-0592:phyxminiproblems-problemphyxmini0592-relativisticseparationsetup}
  Dimensionless relative speed β = v/c
\end{definition}
\begin{proof}
  This is the declared model definition of speed Fraction Of Light.
\end{proof}
\begin{definition}[lorentz Factor]
  \label{def:physics:phyx-mini-0592:phyxminiproblems-problemphyxmini0592-lorentzfactor}
  \lean{PhyXMiniProblems.ProblemPhyXMini0592.lorentzFactor}
  \uses{def:physics:phyx-mini-0592:phyxminiproblems-problemphyxmini0592-relativisticseparationsetup, def:physics:phyx-mini-0592:phyxminiproblems-problemphyxmini0592-speedfractionoflight}
  Physlib's Lorentz factor for the relative speed
\end{definition}
\begin{proof}
  This is the declared model definition of lorentz Factor.
\end{proof}
\begin{definition}[Matches Relativistic Separation Scenario]
  \label{def:physics:phyx-mini-0592:phyxminiproblems-problemphyxmini0592-matchesrelativisticseparationscenario}
  \lean{PhyXMiniProblems.ProblemPhyXMini0592.MatchesRelativisticSeparationScenario}
  \uses{def:physics:phyx-mini-0592:phyxminiproblems-problemphyxmini0592-relativisticseparationsetup}
  Frame roles and boost direction stated or implied by the problem
\end{definition}
\begin{proof}
  This is the declared model definition of Matches Relativistic Separation Scenario.
\end{proof}
\begin{definition}[Matches Supplied Separation Time Graph]
  \label{def:physics:phyx-mini-0592:phyxminiproblems-problemphyxmini0592-matchessuppliedseparationtimegraph}
  \lean{PhyXMiniProblems.ProblemPhyXMini0592.MatchesSuppliedSeparationTimeGraph}
  \uses{def:physics:phyx-mini-0592:phyxminiproblems-problemphyxmini0592-separationinmeters, def:physics:phyx-mini-0592:phyxminiproblems-problemphyxmini0592-timeintervalofnanoseconds, def:physics:phyx-mini-0592:phyxminiproblems-problemphyxmini0592-displayedlineslopemeterspernanosecond, def:physics:phyx-mini-0592:phyxminiproblems-problemphyxmini0592-relativisticseparationsetup}
  Direct transcription of 592.png. The bitmap, rather than its auxiliary caption, determines the nonzero 2 m intercept and the 9 m value at the 10 ns right boundary. The last field states only that the physical graph follows that displayed affine line; it does not mention Δx'
\end{definition}
\begin{proof}
  This is the declared model definition of Matches Supplied Separation Time Graph.
\end{proof}
\begin{definition}[Has Physical Relativistic Parameters]
  \label{def:physics:phyx-mini-0592:phyxminiproblems-problemphyxmini0592-hasphysicalrelativisticparameters}
  \lean{PhyXMiniProblems.ProblemPhyXMini0592.HasPhysicalRelativisticParameters}
  \uses{def:physics:phyx-mini-0592:phyxminiproblems-problemphyxmini0592-separationinmeters, def:physics:phyx-mini-0592:phyxminiproblems-problemphyxmini0592-vacuumlightspeedinmeterspernanosecond, def:physics:phyx-mini-0592:phyxminiproblems-problemphyxmini0592-relativisticseparationsetup, def:physics:phyx-mini-0592:phyxminiproblems-problemphyxmini0592-speedfractionoflight}
  The positivity and subluminality regime required for the Lorentz model. These conditions constrain the physical branch but supply no numerical value for the requested primed separation
\end{definition}
\begin{proof}
  This is the declared model definition of Has Physical Relativistic Parameters.
\end{proof}
\begin{definition}[Satisfies Inverse Lorentz Spatial Transformation]
  \label{def:physics:phyx-mini-0592:phyxminiproblems-problemphyxmini0592-satisfiesinverselorentzspatialtransformation}
  \lean{PhyXMiniProblems.ProblemPhyXMini0592.SatisfiesInverseLorentzSpatialTransformation}
  \uses{def:physics:phyx-mini-0592:phyxminiproblems-problemphyxmini0592-timeinterval, def:physics:phyx-mini-0592:phyxminiproblems-problemphyxmini0592-separationinmeters, def:physics:phyx-mini-0592:phyxminiproblems-problemphyxmini0592-timeintervalinnanoseconds, def:physics:phyx-mini-0592:phyxminiproblems-problemphyxmini0592-speedinmeterspernanosecond, def:physics:phyx-mini-0592:phyxminiproblems-problemphyxmini0592-relativisticseparationsetup, def:physics:phyx-mini-0592:phyxminiproblems-problemphyxmini0592-lorentzfactor}
  Generic inverse spatial Lorentz transformation Δx = γ(β) (Δx' + v Δt'), expressed in the coherent metre/nanosecond unit system used by the graph. It relates independent fields of the setup and contains neither the requested numerical answer nor any answer-choice label
\end{definition}
\begin{proof}
  This is the declared model definition of Satisfies Inverse Lorentz Spatial Transformation.
\end{proof}
\begin{definition}[Answer Choice]
  \label{def:physics:phyx-mini-0592:phyxminiproblems-problemphyxmini0592-answerchoice}
  \lean{PhyXMiniProblems.ProblemPhyXMini0592.AnswerChoice}
  Labels of the four metre-valued answer choices
\end{definition}
\begin{proof}
  This is the declared model definition of Answer Choice.
\end{proof}
\begin{definition}[displayed Separation In Meters]
  \label{def:physics:phyx-mini-0592:phyxminiproblems-problemphyxmini0592-displayedseparationinmeters}
  \lean{PhyXMiniProblems.ProblemPhyXMini0592.displayedSeparationInMeters}
  \uses{def:physics:phyx-mini-0592:phyxminiproblems-problemphyxmini0592-answerchoice}
  Numerical separation printed beside each answer-choice label, in metres
\end{definition}
\begin{proof}
  This is the declared model definition of displayed Separation In Meters.
\end{proof}
\begin{definition}[displayed Half Unit In Last Place Meters]
  \label{def:physics:phyx-mini-0592:phyxminiproblems-problemphyxmini0592-displayedhalfunitinlastplacemeters}
  \lean{PhyXMiniProblems.ProblemPhyXMini0592.displayedHalfUnitInLastPlaceMeters}
  \uses{def:physics:phyx-mini-0592:phyxminiproblems-problemphyxmini0592-answerchoice}
  Half of one unit in the last printed decimal place for each displayed answer
\end{definition}
\begin{proof}
  This is the declared model definition of displayed Half Unit In Last Place Meters.
\end{proof}
\begin{definition}[Matches Displayed Precision]
  \label{def:physics:phyx-mini-0592:phyxminiproblems-problemphyxmini0592-matchesdisplayedprecision}
  \lean{PhyXMiniProblems.ProblemPhyXMini0592.MatchesDisplayedPrecision}
  \uses{def:physics:phyx-mini-0592:phyxminiproblems-problemphyxmini0592-spatialseparation, def:physics:phyx-mini-0592:phyxminiproblems-problemphyxmini0592-separationinmeters, def:physics:phyx-mini-0592:phyxminiproblems-problemphyxmini0592-answerchoice, def:physics:phyx-mini-0592:phyxminiproblems-problemphyxmini0592-displayedseparationinmeters, def:physics:phyx-mini-0592:phyxminiproblems-problemphyxmini0592-displayedhalfunitinlastplacemeters}
  A physical separation agrees with an answer at its displayed precision
\end{definition}
\begin{proof}
  This is the declared model definition of Matches Displayed Precision.
\end{proof}
\begin{definition}[Is Nearest Answer Choice]
  \label{def:physics:phyx-mini-0592:phyxminiproblems-problemphyxmini0592-isnearestanswerchoice}
  \lean{PhyXMiniProblems.ProblemPhyXMini0592.IsNearestAnswerChoice}
  \uses{def:physics:phyx-mini-0592:phyxminiproblems-problemphyxmini0592-spatialseparation, def:physics:phyx-mini-0592:phyxminiproblems-problemphyxmini0592-separationinmeters, def:physics:phyx-mini-0592:phyxminiproblems-problemphyxmini0592-answerchoice, def:physics:phyx-mini-0592:phyxminiproblems-problemphyxmini0592-displayedseparationinmeters}
  The selected choice is at least as close as every displayed alternative
\end{definition}
\begin{proof}
  This is the declared model definition of Is Nearest Answer Choice.
\end{proof}
% --- Archon declaration topology end ---
% --- Archon physics formalization source end ---
